\section{Locality of Gaussian-filtered terms on a finite graph}
\label{sec:hastings_local_projective_form}

% Hastings, arXiv:cond-mat/0508554v3, proj3new.tex, lines 169--190 (setting),
% 266--387 (eqs. lbnd, mid, odef, hpdef, apb, facld and the unique ground
% state), proofs in the appendix at lines 1419--1533.

The approximate local projective form of~\cite{Hastings2006Solving} replaces
each term $H_i$ of a gapped Hamiltonian $H=\sum_kh_k$ by its Gaussian-filtered
operator $\widetilde H_i$ (Definition~\ref{def:hastings_gaussian_filtered_term})
and then by a finite-range approximation $M_i$. Under the finite-range graph
assumptions of Theorem~\ref{thm:qc_graph_lieb_robinson_finset} (each
$h_k=h_k^\dagger$ supported on $X_k$ with $\operatorname{diam}X_k\le R$,
$|X_k|\le v_R$, and $\sum_{k:\,x\in X_k}\|h_k\|\le b_0$), we define $M_i$ by
the normalized partial trace of $\widetilde H_i$ onto the $l$-neighbourhood of
the support of $H_i$, rather than by truncated local evolution
\cite[lines 285--307]{Hastings2006Solving}, and read the Gaussian weight of
eq.~(mid) as $e^{-(t\Delta E)^2/(2q)}$
\cite{gap:hastings05_filtered_term_approximation}. For $q\ge0$ and
$\Delta E\in\mathbb R$, $\widetilde W$ denotes the Gaussian average of
$e^{iHt}We^{-iHt}$ of variance $q/(\Delta E)^2$ when $q>0$ and $\Delta E\ne0$,
and $\widetilde W=W$ when $q=0$ or $\Delta E=0$.

\begin{theorem}[Locality of a filtered term on a graph]
    \label{thm:hastings05_filtered_locality_graph}
    \lean{TNLean.Lattice.norm_commutator_gaussianFilteredTerm_le_graphDistance}
    \leanok
    \uses{thm:hastings_filtered_locality,thm:qc_graph_lieb_robinson_finset}
    Let $W$ act on $S$, $B$ act on $\{y\}$ with $\min_{x\in S}d_G(x,y)\ge l$,
    and $q,\mu,T\ge0$, $\Delta E\in\mathbb R$. With $(C,v)$ from
    Theorem~\ref{thm:qc_graph_lieb_robinson_finset},
    \begin{align}
        \|[\widetilde W,B]\|\le C|S|\,\|W\|\,\|B\|\,e^{vT-\mu l}
        +2\|W\|\,\|B\|\,e^{-(T\Delta E)^2/(2q)}.
        \label{eq:hastings05_filtered_graph_commutator}
    \end{align}
    For $T=c_1l$ this is eq.~(lbnd) of~\cite[lines 266--283]{Hastings2006Solving}.
\end{theorem}
\begin{proof}\leanok
    Combine the $|t|\le T$ commutator bound of
    Theorem~\ref{thm:qc_graph_lieb_robinson_finset} with
    Theorem~\ref{thm:hastings_filtered_locality}.
\end{proof}

\begin{theorem}[The finite-range approximation of a filtered term]
    \label{thm:hastings05_filtered_approximation_graph}
    \lean{TNLean.Lattice.norm_gaussianFilteredTerm_sub_siteExpectation_le_graphNeighborhood,
        TNLean.Lattice.exists_local_approx_gaussianFilteredTerm,
        TNLean.Lattice.isHermitian_gaussianFilteredTerm}
    \leanok
    \uses{thm:hastings_filtered_partial_trace,thm:qc_graph_neighborhood_localization,
        lem:qc_site_expectation_properties,lem:al_site_expectation_order,
        thm:hastings_gaussian_filtered_spectral}
    Assume $|\{y:d_G(x,y)=r\}|\le K_g(r+1)^k$ for all $x$ and $r\ge0$, and let
    $(C,v,T_k)$ be as in Theorem~\ref{thm:qc_graph_neighborhood_localization}.
    For $W$ supported on $S$, $l\in\mathbb N$, $N_l(S)=\{y:d_G(y,S)\le l\}$,
    positive local dimension, $q,T\ge0$, and $\Delta E\in\mathbb R$, the
    operator $M=\mathbb E_{N_l(S)}(\widetilde W)$ acts on $N_l(S)$, satisfies
    $\|M\|\le\|W\|$ (and $M=M^\dagger$ if $W=W^\dagger$), and
    \begin{align}
        \|\widetilde W-M\|\le C|S|^2\,\|W\|\,K_gT_k\,e^{vT-l/2}
        +2\|W\|\,e^{-(T\Delta E)^2/(2q)}.
        \label{eq:hastings05_filtered_graph_approximation}
    \end{align}
    For $W=H_i$ supported in the $R$-ball around $i$ and
    $l=l_{\mathrm{proj}}-R\ge0$, $M_i=M$ acts within radius
    $l_{\mathrm{proj}}$ of $i$ with $\|M_i\|\le J$
    \cite[lines 285--320]{Hastings2006Solving},
    \cite{gap:hastings05_filtered_term_approximation}.
\end{theorem}
\begin{proof}\leanok
    Combine Theorem~\ref{thm:qc_graph_neighborhood_localization} for $|t|\le T$
    with Theorem~\ref{thm:hastings_filtered_partial_trace}.
\end{proof}

\begin{theorem}[Ground vectors are nearly invariant under the approximation]
    \label{thm:hastings05_facld}
    \lean{TNLean.Lattice.norm_sub_starProjection_apply_le_of_gaussianFilteredTerm,
        TNLean.Lattice.norm_apply_sub_inner_smul_le_of_gaussianFilteredTerm}
    \leanok
    \uses{thm:hastings_gaussian_filtered_spectral,
        thm:hastings_filtered_no_ground_component,thm:area_law_ground_eigenspace}
    Let $H=H^\dagger$ have ground energy $E_0$, gap $\Delta E>0$, and ground
    projection $P$, and let $q\ge0$. For any $H\psi=E_0\psi$ and operators
    $W,M$, $\|(1-P)M\psi\|\le(\|W\|e^{-q/2}+\|\widetilde W-M\|)\|\psi\|$. If
    $\operatorname{ran}P=\mathbb C\Omega$ with $\|\Omega\|=1$, then
    \begin{align}
        \|M\Omega-\langle\Omega,M\Omega\rangle\Omega\|
        \le\|W\|e^{-q/2}+\|\widetilde W-M\|
        \label{eq:hastings05_filtered_ground_vector}
    \end{align}
    \cite[eq.~(facld), lines 322--348 and 375--387]{Hastings2006Solving}.
\end{theorem}
\begin{proof}\leanok
    Since $(1-P)M\psi=(1-P)\widetilde W\psi+(1-P)(M-\widetilde W)\psi$,
    Theorem~\ref{thm:hastings_gaussian_filtered_spectral} bounds the first term
    by $\|W\|e^{-q/2}\|\psi\|$ and $\|(1-P)(M-\widetilde W)\psi\|\le\|\widetilde W-M\|\,\|\psi\|$.
    For $\operatorname{ran}P=\mathbb C\Omega$, $P\phi=\langle\Omega,\phi\rangle\Omega$.
\end{proof}

\begin{theorem}[The sum of the approximations]
    \label{thm:hastings05_sum_approximation}
    \lean{TNLean.Lattice.norm_sum_sub_le_of_sum_eq}
    \leanok
    \uses{thm:hastings_gaussian_filtered_spectral}
    If $H=H^\dagger=\sum_iW_i$, then any $M_i$ satisfy
    $\|\sum_iM_i-H\|\le\sum_i\|\widetilde W_i-M_i\|$
    \cite[lines 359--374]{Hastings2006Solving}.
\end{theorem}
\begin{proof}\leanok
    By Theorem~\ref{thm:hastings_gaussian_filtered_spectral},
    $\sum_i\widetilde W_i=H$, so $\sum_iM_i-H=\sum_i(M_i-\widetilde W_i)$.
\end{proof}

\subsection*{Degenerate ground spaces}

% Hastings, arXiv:cond-mat/0508554v3, proj3new.tex, lines 395--491 (degenerate
% ground space), Appendix B at lines 1535--1708 (eqs. lc, blecm, oibu).

For a degenerate ground space with orthogonal projection $P$, ground states are
approximate eigenvectors of $M_i$ in a basis that approximately diagonalizes
the compressions $PM_iP$ \cite[lines 395--491]{Hastings2006Solving}. Throughout
this subsection, $H$ has ground energy $E_0$ and gap $\Delta E>0$, and $q\ge0$.
Separated-operator commutators are bounded via
Theorem~\ref{thm:hastings05_filtered_approximation_graph}
\cite{gap:hastings05_filtered_term_approximation}; the translation-invariant
case and low-energy example are in
Sections~\ref{sec:hastings_translation_supercell}
and~\ref{sec:hastings_low_energy_example}.

\begin{theorem}[Low-energy commutators]
    \label{thm:hastings05_low_energy_commutator}
    \lean{TNLean.Lattice.groundCompression,
        TNLean.Lattice.groundCompression_gaussianFilteredTerm,
        TNLean.Lattice.norm_one_sub_starProjection_mul_gaussianFilteredTerm_mul_le,
        TNLean.Lattice.norm_starProjection_mul_gaussianFilteredTerm_mul_one_sub_le,
        TNLean.Lattice.norm_groundCompression_commutator_le}
    \leanok
    \uses{thm:hastings_gaussian_filtered_spectral}
    For all operators $W,B$, $P\widetilde WP=PWP$,
    $\|(1-P)\widetilde WP\|\le\|W\|e^{-q/2}$,
    $\|P\widetilde W(1-P)\|\le\|W\|e^{-q/2}$, and
    \begin{align}
        \|[PWP,PBP]\|\le\|[\widetilde W,B]\|+2\|W\|\,\|B\|\,e^{-q/2}
        \label{eq:hastings05_low_energy_bound}
    \end{align}
    \cite[lines 1562--1593]{Hastings2006Solving}.
\end{theorem}
\begin{proof}\leanok
    Theorem~\ref{thm:hastings_gaussian_filtered_spectral} gives
    $P\widetilde WP=PWP$, $\|(1-P)\widetilde WP\|\le\|W\|e^{-q/2}$, and its
    adjoint bound since $(\widetilde W)^\dagger=\widetilde{W^\dagger}$. With
    $Q=1-P$,
    \begin{align}
        [P\widetilde WP,PBP]=P[\widetilde W,B]P-P\widetilde WQBP
        +PBQ\widetilde WP,
        \label{eq:hastings05_compressed_commutator_split}
    \end{align}
    which yields (\ref{eq:hastings05_low_energy_bound}).
\end{proof}

\begin{theorem}[The commutator with a separated operator]
    \label{thm:hastings05_lc}
    \lean{TNLean.Lattice.norm_commutator_gaussianFilteredTerm_le_of_disjoint}
    \leanok
    \uses{thm:hastings05_filtered_approximation_graph}
    In the setting of Theorem~\ref{thm:hastings05_filtered_approximation_graph},
    if $W$ acts on $S$, $B$ acts on a set disjoint from $N_l(S)$, and $q>0$,
    then for all $T\ge0$,
    \begin{align}
        \|[\widetilde W,B]\|\le2\|B\|\Bigl(C|S|^2\,\|W\|\,K_gT_k\,e^{vT-l/2}
        +2\|W\|\,e^{-(T\Delta E)^2/(2q)}\Bigr)
        \label{eq:hastings05_separated_commutator}
    \end{align}
    (with $e^{-(T\Delta E)^2/(2q)}$ replaced by $1$ at $q=0$)
    \cite[lines 1548--1561]{Hastings2006Solving}.
\end{theorem}
\begin{proof}\leanok
    Since $\mathbb E_{N_l(S)}(\widetilde W)$ commutes with $B$,
    $\|[\widetilde W,B]\|=\|[\widetilde W-\mathbb E_{N_l(S)}(\widetilde W),B]\|
        \le2\|B\|\,\|\widetilde W-\mathbb E_{N_l(S)}(\widetilde W)\|$; apply
    (\ref{eq:hastings05_filtered_graph_approximation}).
\end{proof}

\begin{theorem}[Low-energy commutators of separated operators]
    \label{thm:hastings05_blecm}
    \lean{TNLean.Lattice.norm_groundCompression_commutator_le_graphNeighborhood}
    \leanok
    \uses{thm:hastings05_low_energy_commutator,thm:hastings05_lc}
    Under the assumptions of Theorem~\ref{thm:hastings05_lc},
    \begin{align}
        \|[PWP,PBP]\|\le2\|B\|\Bigl(C|S|^2\,\|W\|\,K_gT_k\,e^{vT-l/2}
        +2\|W\|\,e^{-(T\Delta E)^2/(2q)}\Bigr)+2\|W\|\,\|B\|\,e^{-q/2}
        \label{eq:hastings05_separated_compression}
    \end{align}
    (with $e^{-(T\Delta E)^2/(2q)}$ replaced by $1$ at $q=0$). For $l>0$,
    $0<vc_1<1/2$, $T=c_1l$, and $q=c_1l\Delta E$, the right-hand side of
    (\ref{eq:hastings05_separated_compression}) is
    \begin{align}
        2C|S|^2\|W\|\,\|B\|K_gT_k e^{-(1/2-vc_1)l}
        +6\|W\|\,\|B\|e^{-c_1l\Delta E/2}
        \label{eq:hastings05_exponential_commutator_decay}
    \end{align}
    \cite[lines 1562--1601]{Hastings2006Solving}.
\end{theorem}
\begin{proof}\leanok
    Combine (\ref{eq:hastings05_low_energy_bound}) and
    (\ref{eq:hastings05_separated_commutator}).
\end{proof}

\begin{theorem}[Off-diagonal elements from a commutator]
    \label{thm:hastings05_offdiagonal_commutator}
    \lean{TNLean.Lattice.abs_sub_mul_norm_inner_le_norm_commutator,
        TNLean.Lattice.norm_inner_le_norm_commutator_div}
    \leanok
    If $A=A^\dagger$, $Av=rv$, and $Aw=sw$ with $r\ne s\in\mathbb R$, then for
    every $B$,
    \begin{align}
        |\langle v,Bw\rangle|\le\frac{\|[A,B]\|\,\|v\|\,\|w\|}{|r-s|}
        \label{eq:hastings05_offdiagonal_bound}
    \end{align}
    \cite[lines 1678--1690]{Hastings2006Solving}.
\end{theorem}
\begin{proof}\leanok
    Use $(r-s)\langle v,Bw\rangle=\langle v,[A,B]w\rangle$ and Cauchy--Schwarz.
\end{proof}

\begin{theorem}[A basis of a two-dimensional ground space]
    \label{thm:hastings05_two_dim_basis}
    \lean{TNLean.Lattice.offDiagonalPart,
        TNLean.Lattice.offDiagonalPart_apply,
        TNLean.Lattice.norm_offDiagonalPart_le,
        TNLean.Lattice.norm_offDiagonalPart_le_norm_commutator_div,
        TNLean.Lattice.sum_norm_offDiagonalPart_le}
    \leanok
    \uses{thm:hastings05_offdiagonal_commutator}
    For orthonormal $v,w$ and
    $o_B=\langle v,Bw\rangle\lvert v\rangle\langle w\rvert
        +\langle w,Bv\rangle\lvert w\rangle\langle v\rvert$, one has
    $\|o_B\|\le\max(|\langle v,Bw\rangle|,|\langle w,Bv\rangle|)$. If
    $A=A^\dagger$ satisfies $Av=rv$ and $Aw=sw$ with $r\ne s$, then
    \begin{align}
        \|o_B\|\le\frac{\|[A,B]\|}{|r-s|},
        \qquad
        \sum_i\|o_{B_i}\|\le\frac{\sum_i\|[A,B_i]\|}{|r-s|}
        \label{eq:hastings05_two_dim_basis}
    \end{align}
    \cite[lines 1678--1694]{Hastings2006Solving},
    \cite{gap:hastings05_two_dim_ground_basis}.
\end{theorem}
\begin{proof}\leanok
    By Bessel's inequality,
    $\|o_Bx\|^2=|\langle v,Bw\rangle|^2|\langle w,x\rangle|^2
        +|\langle w,Bv\rangle|^2|\langle v,x\rangle|^2
        \le\max(|\langle v,Bw\rangle|,|\langle w,Bv\rangle|)^2\|x\|^2$; apply
    (\ref{eq:hastings05_offdiagonal_bound}).
\end{proof}

\begin{theorem}[Ground states as approximate eigenvectors up to the off-diagonal part]
    \label{thm:hastings05_low_up_to_offdiagonal}
    \lean{TNLean.Lattice.norm_apply_sub_inner_smul_le_add_offDiagonalPart}
    \leanok
    \uses{thm:hastings05_two_dim_basis}
    For orthonormal $v,w$, $P=\lvert v\rangle\langle v\rvert+\lvert w\rangle\langle w\rvert$,
    and any $M$,
    \begin{align}
        \|Mv-\langle v,Mv\rangle v\|\le\|(1-P)Mv\|+\|o_{PMP}\|
        \label{eq:hastings05_low_up_to_offdiagonal}
    \end{align}
    \cite[eq.~(low), lines 476--491]{Hastings2006Solving}.
\end{theorem}
\begin{proof}\leanok
    Since $PMv-\langle v,Mv\rangle v=\langle w,PMPv\rangle w=o_{PMP}v$, the
    triangle inequality gives (\ref{eq:hastings05_low_up_to_offdiagonal}).
\end{proof}

For $\dim\operatorname{ran}P=n>2$, we treat simple spectrum of $PM_jP$ first
and then the iterative splitting of repeated eigenvalues
\cite[lines 1695--1705]{Hastings2006Solving}.

\begin{theorem}[The off-diagonal part for a compression with simple spectrum]
    \label{thm:hastings05_simple_spectrum_basis}
    \lean{TNLean.Lattice.shiftPart,
        TNLean.Lattice.shiftPart_apply,
        TNLean.Lattice.norm_shiftPart_le,
        TNLean.Lattice.offDiagonalPartFamily,
        TNLean.Lattice.norm_offDiagonalPartFamily_le,
        TNLean.Lattice.norm_offDiagonalPartFamily_le_norm_commutator_div,
        TNLean.Lattice.sum_norm_offDiagonalPartFamily_le}
    \leanok
    \uses{thm:hastings05_offdiagonal_commutator}
    Let $v_0,\dots,v_{n-1}$ be orthonormal and
    $o_B=\sum_{a\ne b}\langle v_a,Bv_b\rangle\lvert v_a\rangle\langle v_b\rvert$.
    If $|\langle v_a,Bv_b\rangle|\le m$ for all $a\ne b$, then
    $\|o_B\|\le(n-1)m$. If $A=A^\dagger$ satisfies $Av_a=r_av_a$ with
    $|r_a-r_b|\ge\delta>0$ for all $a\ne b$, then
    \begin{align}
        \|o_B\|\le(n-1)\frac{\|[A,B]\|}{\delta},
        \qquad
        \sum_i\|o_{B_i}\|\le(n-1)\frac{\sum_i\|[A,B_i]\|}{\delta}
        \label{eq:hastings05_simple_spectrum_basis}
    \end{align}
    \cite[lines 1678--1705]{Hastings2006Solving},
    \cite{gap:hastings05_two_dim_ground_basis}.
\end{theorem}
\begin{proof}\leanok
    Decompose $o_B=\sum_{k=1}^{n-1}S_k$ with
    $S_k=\sum_a\langle v_a,Bv_{a+k}\rangle\lvert v_a\rangle\langle v_{a+k}\rvert$
    (indices mod $n$). Bessel's inequality gives $\|S_k\|\le m$, and
    (\ref{eq:hastings05_offdiagonal_bound}) bounds each entry by
    $\|[A,B]\|/\delta$.
\end{proof}

\begin{theorem}[Ground states as approximate eigenvectors, any dimension]
    \label{thm:hastings05_low_up_to_offdiagonal_family}
    \lean{TNLean.Lattice.offDiagonalPartFamily_apply_self,
        TNLean.Lattice.norm_apply_sub_inner_smul_le_add_offDiagonalPartFamily}
    \leanok
    \uses{thm:hastings05_simple_spectrum_basis}
    For an orthonormal basis $v_0,\dots,v_{n-1}$ of $\operatorname{ran}P$, any
    $M$, and any $a$,
    \begin{align}
        \|Mv_a-\langle v_a,Mv_a\rangle v_a\|\le\|(1-P)Mv_a\|+\|o_{PMP}\|
        \label{eq:hastings05_low_up_to_offdiagonal_family}
    \end{align}
    \cite[eq.~(low), lines 476--491]{Hastings2006Solving}.
\end{theorem}
\begin{proof}\leanok
    Use $PMv_a-\langle v_a,Mv_a\rangle v_a
        =\sum_{b\ne a}\langle v_b,PMPv_a\rangle v_b=o_{PMP}v_a$.
\end{proof}

For an orthonormal basis $v_0,\dots,v_{n-1}$ of $\operatorname{ran}P$ and real
labels $r_\ell(a)$ ($\ell\ge0$), let the level-$\ell$ class of $a$ be
$C_\ell(a)=\{c:\forall m<\ell,\ r_m(c)=r_m(a)\}$, assume
$\langle v_c,A_\ell v_d\rangle=r_\ell(c)\delta_{cd}$ for $c,d\in C_\ell(a)$ where
$A_\ell=PM_{j_\ell}P$, and let $\delta_\ell>0$ be the minimum of
$|r_\ell(a)-r_\ell(b)|>0$ over $b\in C_\ell(a)$ ($\delta_\ell=1$ if no such
pair exists) \cite[lines 1695--1705]{Hastings2006Solving}.

\begin{theorem}[One level of the iteration]
    \label{thm:hastings05_basis_iteration_step}
    \lean{TNLean.Lattice.inner_commutator_apply_eq_sum,
        TNLean.Lattice.abs_sub_mul_norm_inner_le}
    \leanok
    If $v_0,\dots,v_{n-1}$ are orthonormal, $A,B$ preserve their span, and
    $\langle v_c,Av_d\rangle=r_c\delta_{cd}$ for $c,d\in S$, then for $a,b\in S$,
    \begin{align}
        |r_a-r_b|\,|\langle v_a,Bv_b\rangle|
        \le\|[A,B]\|+\sum_{c\notin S}\bigl(|\langle v_a,Av_c\rangle|\,
        |\langle v_c,Bv_b\rangle|+|\langle v_a,Bv_c\rangle|\,
        |\langle v_c,Av_b\rangle|\bigr).
        \label{eq:hastings05_basis_iteration_step}
    \end{align}
\end{theorem}
\begin{proof}\leanok
    Expand $\langle v_a,[A,B]v_b\rangle
        =\sum_c(\langle v_a,Av_c\rangle\langle v_c,Bv_b\rangle
        -\langle v_a,Bv_c\rangle\langle v_c,Av_b\rangle)$; the terms $c\in S$
    contribute $(r_a-r_b)\langle v_a,Bv_b\rangle$.
\end{proof}

\begin{definition}[The levels of the iteration]
    \label{def:hastings05_basis_iteration_levels}
    \lean{TNLean.Lattice.labelClass,
        TNLean.Lattice.levelPairs,
        TNLean.Lattice.levelGap}
    \leanok
    The level-$\ell$ class $C_\ell(a)$, the set of pairs separated first at
    level $\ell$, and the minimal label gap $\delta_\ell$ (with $\delta_\ell=1$
    when empty).
\end{definition}

\begin{theorem}[The bound at the first separating level]
    \label{thm:hastings05_basis_iteration}
    \lean{TNLean.Lattice.levelGap_pos,
        TNLean.Lattice.levelGap_le,
        TNLean.Lattice.norm_inner_le_of_mem_labelClass}
    \leanok
    \uses{thm:hastings05_basis_iteration_step,
        def:hastings05_basis_iteration_levels}
    We have $\delta_\ell>0$, and if $b\in C_\ell(a)$ with
    $r_\ell(a)\ne r_\ell(b)$, then $|\langle v_a,Bv_b\rangle|$ is bounded by
    $\delta_\ell^{-1}$ times the right-hand side of
    (\ref{eq:hastings05_basis_iteration_step}) for $A=A_\ell$ and
    $S=C_\ell(a)$~\cite{gap:hastings05_two_dim_ground_basis}.
\end{theorem}
\begin{proof}\leanok
    Divide (\ref{eq:hastings05_basis_iteration_step}) by
    $|r_\ell(a)-r_\ell(b)|\ge\delta_\ell>0$.
\end{proof}

\begin{theorem}[Two levels]
    \label{thm:hastings05_basis_two_levels}
    \lean{TNLean.Lattice.norm_inner_le_of_two_levels,
        TNLean.Lattice.norm_offDiagonalPartFamily_le_of_two_levels}
    \leanok
    \uses{thm:hastings05_basis_iteration, thm:hastings05_simple_spectrum_basis}
    If every pair $a\ne b$ is separated at level $0$ or $1$, then for all
    $a\ne b$,
    \begin{align}
        |\langle v_a,Bv_b\rangle| & \le
        \max\Bigl(\frac{\|[A_0,B]\|}{\delta_0},\
        \frac{1}{\delta_1}\Bigl(\|[A_1,B]\|
        +2n\frac{\|[A_0,A_1]\|\,\|[A_0,B]\|}{\delta_0^2}\Bigr)\Bigr),
        \label{eq:hastings05_basis_two_levels}
    \end{align}
    and $\|o_B\|$ is at most $n-1$ times the right-hand side of
    (\ref{eq:hastings05_basis_two_levels})
    \cite[lines 1695--1705]{Hastings2006Solving}.
\end{theorem}
\begin{proof}\leanok
    Apply Theorem~\ref{thm:hastings05_basis_iteration} at $\ell=0$ (where
    $S^c=\emptyset$) and at $\ell=1$ (where each $c\notin C_1(a)$ is separated
    from $a,b$ at level $0$), then Theorem~\ref{thm:hastings05_simple_spectrum_basis}.
\end{proof}

\begin{theorem}[Every level]
    \label{thm:hastings05_basis_levels}
    \lean{TNLean.Lattice.exists_lt_mem_labelClass_ne,
        TNLean.Lattice.sum_cross_le,
        TNLean.Lattice.levelBoundA,
        TNLean.Lattice.levelBoundB,
        TNLean.Lattice.norm_inner_A_le_levelBoundA,
        TNLean.Lattice.norm_inner_le_levelBoundB,
        TNLean.Lattice.norm_offDiagonalPartFamily_le_levelBoundB}
    \leanok
    \uses{thm:hastings05_basis_iteration, thm:hastings05_simple_spectrum_basis}
    Assume $\|[A_\ell,A_m]\|\le K$ and $\|[A_\ell,B]\|\le K_B$ for all
    $\ell,m$, and define $\alpha_0=K/\delta_0$, $\beta_0=K_B/\delta_0$,
    \begin{align}
        \alpha_{\ell+1} & =\max\Bigl(\alpha_\ell,\frac{K+2n\alpha_\ell^2}{\delta_{\ell+1}}\Bigr),          &
        \beta_{\ell+1}  & =\max\Bigl(\beta_\ell,\frac{K_B+2n\alpha_\ell\beta_\ell}{\delta_{\ell+1}}\Bigr).
        \label{eq:hastings05_basis_levels_recurrence}
    \end{align}
    If $(a,b)$ is separated first at level $\ell$, then
    $|\langle v_a,Bv_b\rangle|\le\beta_\ell$; if every pair $a\ne b$ is
    separated below level $L$, then $\|o_B\|\le(n-1)\beta_{L-1}$
    \cite[lines 1695--1708]{Hastings2006Solving}.
\end{theorem}
\begin{proof}\leanok
    Induct on $\ell$ using Theorem~\ref{thm:hastings05_basis_iteration} (first
    for $B=A_m$ to bound $|\langle v_a,A_mv_c\rangle|\le\alpha_\ell$, then for
    $B$), and conclude with Theorem~\ref{thm:hastings05_simple_spectrum_basis}.
\end{proof}

\begin{theorem}[The low-energy commutators decay exponentially]
    \label{thm:hastings05_blecm_exp}
    \lean{TNLean.Lattice.lrVelocity,
        TNLean.Lattice.commutatorRate,
        TNLean.Lattice.commutatorConst,
        TNLean.Lattice.norm_groundCompression_le,
        TNLean.Lattice.norm_groundCompression_commutator_le_two_mul,
        TNLean.Lattice.norm_groundCompression_commutator_le_exp,
        TNLean.Lattice.disjoint_biUnion_graphBall_of_lt_edist,
        TNLean.Lattice.norm_groundCompression_commutator_le_exp_edist}
    \leanok
    \uses{thm:hastings05_blecm, thm:hastings05_low_energy_commutator}
    In the setting of Theorem~\ref{thm:hastings05_blecm}, set
    $v=2v_R|b_0|e^{2R}+1$ and $\gamma=\min(1/4,\Delta E/(8v))$. There is an
    explicit constant $C$ such that $\|[Ph_jP,Ph_iP]\|\le C\|h_j\|\,\|h_i\|e^{-\gamma l}$
    whenever $X_i\cap N_l(X_j)=\emptyset$, and
    $\|[Ph_jP,Ph_iP]\|\le Ce^{\gamma(2R+1)}\|h_j\|\,\|h_i\|e^{-\gamma d_G(a_j,a_i)}$
    for anchors $a_j\in X_j$, $a_i\in X_i$.
\end{theorem}
\begin{proof}\leanok
    Apply Theorem~\ref{thm:hastings05_blecm} with $T=l/(4v)$ and $q=T\Delta E$,
    so $vT-l/2=-l/4$ and $q/2=l\Delta E/(8v)$. If $d_G(a_j,a_i)>2R+l$, then
    $X_i\cap N_l(X_j)=\emptyset$; otherwise
    $\|[Ph_jP,Ph_iP]\|\le2\|h_j\|\,\|h_i\|$.
\end{proof}

\begin{theorem}[The sum over all terms, eq.~(oibu)]
    \label{thm:hastings05_oibu}
    \lean{TNLean.Lattice.add_one_pow_mul_exp_neg_le,
        TNLean.Lattice.shellSumConst,
        TNLean.Lattice.sum_add_one_pow_mul_exp_neg_le,
        TNLean.Lattice.sum_exp_neg_edist_le,
        TNLean.Lattice.isSelfAdjoint_groundCompression,
        TNLean.Lattice.sum_norm_groundCompression_commutator_le,
        TNLean.Lattice.sum_norm_offDiagonalPartFamily_le_oibu}
    \leanok
    \uses{thm:hastings05_blecm_exp, thm:hastings05_simple_spectrum_basis}
    If $|\{y:d_G(x,y)=r\}|\le K_g(r+1)^k$, then for every term $j$,
    \begin{align}
        \sum_i\|[Ph_jP,Ph_iP]\|
         & \le Ce^{\gamma(2R+1)}\,|b_0|\,K_g\,S_{\gamma,k}\,\|h_j\|,
        \qquad
        S_{\gamma,k}=\frac{k!\,(2/\gamma)^ke^{\gamma/2}}{1-e^{-\gamma/2}},
        \label{eq:hastings05_oibu}
    \end{align}
    and in an eigenbasis of $PM_jP=Ph_jP$ with eigenvalue gap $\delta>0$,
    $\sum_i\|o_{Ph_iP}\|$ is at most $(n-1)/\delta$ times (\ref{eq:hastings05_oibu}),
    uniformly in system size \cite[eq.~(oibu), lines 1689--1694]{Hastings2006Solving},
    \cite{gap:hastings05_two_dim_ground_basis}.
\end{theorem}
\begin{proof}\leanok
    Group $h_i$ by anchor $x$: $\sum_{a_i=x}\|h_i\|\le b_0$, and
    $\sum_xe^{-\gamma d_G(a_j,x)}\le K_g\sum_{r\ge0}(r+1)^ke^{-\gamma r}\le K_gS_{\gamma,k}$.
    Apply Theorem~\ref{thm:hastings05_simple_spectrum_basis}.
\end{proof}

\begin{theorem}[Eq.~(oibu) for a degenerate ground space]
    \label{thm:hastings05_oibu_levels}
    \lean{TNLean.Lattice.levelBoundB_eq_mul,
        TNLean.Lattice.sum_norm_offDiagonalPartFamily_le_levelBoundB,
        TNLean.Lattice.groundCompression_apply_eq_sum,
        TNLean.Lattice.sum_norm_offDiagonalPartFamily_le_oibu_levels}
    \leanok
    \uses{thm:hastings05_basis_levels, thm:hastings05_oibu}
    In the setting of Theorem~\ref{thm:hastings05_basis_levels} with
    $A_\ell=Ph_{j_\ell}P$ ($\ell<L$) separating all basis pairs and
    $\|h_i\|\le J$, let $\beta_{L-1}$ be defined by
    (\ref{eq:hastings05_basis_levels_recurrence}) at $K=2J^2$ and $K_B=1$. Then
    \begin{align}
        \sum_i\|o_{Ph_iP}\| & \le(n-1)\,\beta_{L-1}\sum_{\ell<L}Ce^{\gamma(2R+1)}\,|b_0|\,K_g\,
        S_{\gamma,k}\,\|h_{j_\ell}\|
        \label{eq:hastings05_oibu_levels}
    \end{align}
    \cite[lines 1689--1708]{Hastings2006Solving}.
\end{theorem}
\begin{proof}\leanok
    By homogeneity of $\beta_{L-1}$ in $K_B$, apply
    Theorem~\ref{thm:hastings05_basis_levels} with
    $K_B(i)=\sum_{\ell<L}\|[A_\ell,Ph_iP]\|$ and sum over $i$ via
    Theorem~\ref{thm:hastings05_oibu}.
\end{proof}
